\documentclass[11pt,a4paper]{article}
\usepackage[utf8]{inputenc}
\usepackage[T1]{fontenc}
\usepackage{geometry}
\usepackage{amsmath,amssymb,amsfonts}
\usepackage{booktabs}
\usepackage{array}
\usepackage{tabularx}
\usepackage{enumitem}
\usepackage{hyperref}
\usepackage{verbatim}
\usepackage{float}
\usepackage{xcolor}

\geometry{margin=1in}
\hypersetup{hidelinks}
\setlist[itemize]{leftmargin=1.3em,itemsep=0pt,topsep=0.25em,parsep=0pt,partopsep=0pt}
\setlist[enumerate]{leftmargin=1.5em,itemsep=0.25em,topsep=0.35em,parsep=0pt,partopsep=0pt}

\newcolumntype{Y}{>{\raggedright\arraybackslash}X}
\newcommand{\criteriaheading}[1]{\par\smallskip\noindent\textbf{#1}\par}
\newcommand{\scoreitem}[3]{\item #1 -- \textit{#2:} #3}
\newcommand{\templatedesc}[2]{\noindent\textbf{#1} \textit{#2}\par}
\newcommand{\boxfield}[1]{\underline{\hspace{#1}}}
\newcommand{\decisionsummarybox}{%
  \noindent
  \begingroup
  \setlength{\fboxsep}{0pt}%
  \setlength{\fboxrule}{1pt}%
  \fcolorbox{gray!65}{white}{%
    \begin{minipage}{0.96\textwidth}
      \colorbox{gray!65}{%
        \parbox{\dimexpr\linewidth-2\fboxsep\relax}{%
          \textcolor{white}{\strut\textbf{Platform Decision Summary}}%
        }%
      }

      \vspace{0.55em}
      {\small
      \begin{tabularx}{\linewidth}{@{}>{\raggedright\arraybackslash}p{0.33\linewidth}X>{\raggedright\arraybackslash}p{0.23\linewidth}X@{}}
        \textbf{Platform:} & \boxfield{5.1cm} & & \\
        \textbf{Use-Case Template:} & \boxfield{3.8cm} & \textbf{Weighting Version:} & \boxfield{2.6cm} \\
        \textbf{Normalized Score:} & \boxfield{1.4cm}\ \% & \textbf{Raw Weighted:} & \boxfield{1.4cm}\ /\ \boxfield{1.4cm} \\
        \textbf{Go/No-Go Threshold:} & \boxfield{1.4cm}\ \% & \textbf{Decision:} & \boxfield{2.8cm} \\
        \textbf{Knock-Out Status:} & \multicolumn{3}{l}{\boxfield{9.1cm}} \\
      \end{tabularx}

      \vspace{0.45em}
      \textbf{Top 3 Strengths:}\par
      \hspace*{3.8cm}1.\ \boxfield{7.1cm}\par
      \hspace*{3.8cm}2.\ \boxfield{7.1cm}\par
      \hspace*{3.8cm}3.\ \boxfield{7.1cm}\par

      \vspace{0.35em}
      \textbf{Critical Gaps (if any):}\par
      \boxfield{12.5cm}\par

      \vspace{0.35em}
      \textbf{Mitigations / Next Steps:}\par
      \boxfield{12.5cm}\par

      \vspace{0.35em}
      \textbf{Reviewer / Date:}\ \boxfield{3.1cm}\ /\ \boxfield{2.7cm}\par
      }%
    \end{minipage}%
  }%
  \endgroup
}

\title{SDR Hardware Comparison}
\date{May 7, 2026}

\begin{document}

\maketitle

\section*{1 Start Workflow}

This framework provides a repeatable procedure for evaluating compute hosts and SDR RF front-ends for spectrum sensing and radio-based processing. The goal is to select a cost-effective, compact, modular, and scalable sensing node while reducing vendor lock-in.

A candidate configuration should be evaluated as an integrated system:

\[
\text{Sensing Node} \equiv \text{Compute Host} \equiv \text{SDR RF Front-End} \equiv \text{Timing and Synchronization} \equiv \text{Deployment Infrastructure}
\]

Before scoring, verify mandatory requirements such as frequency coverage, instantaneous bandwidth, sustained data throughput, software support, timing capability, power budget, thermal behavior, and calibration feasibility.

The following workflow guides the selection process:

\begin{enumerate}
  \item Define the use case. Specify target bands, bandwidth, signal types, real-time needs, ML requirements, deployment conditions, and budget.
  \item Check mandatory requirements. Reject or flag candidates that fail critical RF, compute, timing, power, thermal, software, or calibration requirements.
  \item Shortlist candidates. Select three to five feasible combinations of compute host, SDR front-end, timing reference, storage/network interface, and power solution.
  \item Evaluate compute-host suitability. Assess CPU/GPU capability, memory, I/O bandwidth, storage, container support, software ecosystem, security, manageability, and sustained operation.
  \item Evaluate SDR front-end suitability. Assess frequency range, instantaneous bandwidth, ADC/DAC resolution, dynamic range, frequency stability, clocking, timestamping, gain control, DC offset/IQ-imbalance behavior, calibration support, and SDR software compatibility and support.
  \item Apply the scoring rubric. Score each applicable criterion from 1 to 5 and multiply it by its use-case weight. Mark non-applicable criteria as N/A.
  \item Assess deployment strategy. Consider single-node versus multi-node readiness, enclosure, cooling, power, networking, data logging, remote access, and maintenance.
  \item Validate with a pilot. Run representative GNU Radio or Python workflows and verify throughput, latency, packet loss, buffer stability, CPU/GPU load, storage bandwidth, temperature, and long-term stability.
  \item Document the decision. Record scores, selected configuration, rejected alternatives, risks, mitigations, calibration assumptions, lock-in risk, openness, and total cost of ownership.
\end{enumerate}

\section*{2 Evaluation Criteria}

\subsection*{2.1 Scoring Rubric}

Apply a 1--5 raw score to each applicable criterion and multiply it by the weight assigned to the selected deployment profile:

\[
\text{Weighted Score}_i = \text{Raw Score}_i \times \text{Weight}_i.
\]

\begin{table}[H]
  \centering
  \small
  \begin{tabularx}{\textwidth}{@{}>{\raggedright\arraybackslash}p{0.36\textwidth}*{4}{>{\centering\arraybackslash}X}@{}}
    \toprule
    Criterion & Template 1: SIGINT & Template 2: Field & Template 3: Lab & Template 4: IoT \\
    \midrule
    Compute capability & 2$\times$ & 2$\times$ & 3$\times$ & 2$\times$ \\
    SDR software support & 2$\times$ & 2$\times$ & 3$\times$ & 2$\times$ \\
    I/O and data throughput & 2$\times$ & 2$\times$ & 3$\times$ & 2$\times$ \\
    Power and thermal behavior & 1$\times$ & 3$\times$ & 1$\times$ & 2$\times$ \\
    Form factor & 1$\times$ & 3$\times$ & 1$\times$ & 2$\times$ \\
    Software ecosystem & 2$\times$ & 2$\times$ & 3$\times$ & 3$\times$ \\
    Expandability & 1$\times$ & 1$\times$ & 3$\times$ & 2$\times$ \\
    Openness / lock-in risk & 2$\times$ & 1$\times$ & 3$\times$ & 2$\times$ \\
    Reliability & 3$\times$ & 3$\times$ & 1$\times$ & 3$\times$ \\
    Total cost of ownership & 2$\times$ & 2$\times$ & 2$\times$ & 3$\times$ \\
    Multi-node readiness & 3$\times$ & 1$\times$ & 1$\times$ & 3$\times$ \\
    Security and manageability & 3$\times$ & 1$\times$ & 1$\times$ & 3$\times$ \\
    RF measurement fidelity & 3$\times$ & 2$\times$ & 2$\times$ & 2$\times$ \\
    Timing and synchronization & 3$\times$ & 2$\times$ & 2$\times$ & 2$\times$ \\
    FPGA/HDL access & 2$\times$ & 1$\times$ & 3$\times$ & 1$\times$ \\
    Regulatory and calibration support & 2$\times$ & 2$\times$ & 1$\times$ & 3$\times$ \\
    \midrule
    Final Weighted Score & & & & \\
    \bottomrule
  \end{tabularx}
  \caption{Use-case weighting templates for candidate sensing-node configurations}
\end{table}

\templatedesc{Template 1: Spectrum Monitoring \& SIGINT}{Fixed infrastructure, continuous operation, multi-site coherent sensing.}
\templatedesc{Template 2: Field Research \& Portable Operations}{Battery-powered, ruggedized, single-operator field collection.}
\templatedesc{Template 3: Laboratory Research \& Algorithm Development}{Benchtop experimentation, flexible DSP, and FPGA development.}
\templatedesc{Template 4: Production IoT/Wireless Monitoring}{Cost-optimized fleet deployment, long lifecycle, and remote management.}

\subsubsection*{Assessment Summary}

\begin{itemize}
  \item \textit{Pros:} [User fills in narrative pros here.]
  \item \textit{Cons/Caveats:} [User fills in narrative cons here.]
\end{itemize}

\subsection*{2.2 Score Normalization \& Decision Logic}

\textit{Normalization Formula.} To enable cross-platform comparison, raw weighted scores are converted to a percentage:

\begin{equation}
\label{eq:normalized-score}
\text{Normalized Score (\%)} =
\frac{\sum_{i=1}^{N} \left(\text{Raw Score}_i \times \text{Weight}_i\right)}
{\sum_{i=1}^{N} \left(5 \times \text{Weight}_i\right)} \times 100.
\end{equation}

Here, $N$ is the number of applicable criteria, excluding those marked N/A.

\noindent\textit{N/A Handling:} For criteria marked N/A, exclude both the raw weighted score from the numerator and the corresponding maximum weighted score from the denominator in Eq.~\eqref{eq:normalized-score}.

\subsection*{2.3 Decision Thresholds}

These thresholds apply only to candidates that have already passed the mandatory requirements gate.

\begin{table}[H]
  \centering
  \small
  \begin{tabular}{@{}lll@{}}
    \toprule
    Normalized Score & Recommendation & Guidance \\
    \midrule
    $\geq 85\%$ & \textbf{Strong Accept} & Proceed to pilot; minimal risk \\
    $70\%$--$84\%$ & \textbf{Conditional Accept} & Proceed with documented mitigations \\
    $55\%$--$69\%$ & \textbf{Re-evaluate} & Address low-scoring criteria or adjust weights \\
    $< 55\%$ & \textbf{Reject} & Platform unlikely to meet requirements \\
    \bottomrule
  \end{tabular}
  \caption{Decision thresholds based on normalized score}
  \label{tab:decision-thresholds}
\end{table}

\section*{3 Implementation Notes}

\begin{itemize}
  \item \textit{Host interfaces.} Prioritize USB 3.x, PCIe, and Gigabit Ethernet, and verify sustained throughput at the target sample rate without buffer overflows or dropped packets.
  \item \textit{Software portability.} Use containerization when possible, but verify SDR-driver access and GPU/accelerator visibility inside the container.
  \item \textit{Power and thermal stability.} For field or continuous deployments, use a robust power supply and adequate cooling. Thermal drift and throttling can affect both processing stability and RF phase/frequency stability.
  \item \textit{Open ecosystem.} Prefer maintained drivers, open APIs, and active community support to reduce vendor lock-in and improve long-term maintainability.
  \item \textit{Machine-learning workloads.} If ML inference is required, verify accelerator availability, memory capacity, model-runtime support, and compatibility with the selected SDR pipeline.
  \item \textit{RF impairments.} Evaluate IQ imbalance, DC offset, gain calibration, and the availability of driver- or software-level compensation hooks.
  \item \textit{FPGA/HDL development.} If FPGA customization is required, account for HDL toolchain availability, license cost, supported devices, and integration with the SDR software stack.
\end{itemize}

\clearpage
\appendix

\section*{Appendix}

\subsection*{Pilot Validation}

\begin{enumerate}
  \item \textbf{Throughput \& Stability Test}
  \begin{itemize}
    \item Run continuous acquisition at target sample rate for 4+ hours.
    \item Monitor for overflows, underruns, dropped samples.
    \item \textit{Pass criterion:} $<0.01\%$ sample loss, no crashes.
    \item \textit{Tools:} \texttt{uhd\_rx\_cfile}, \texttt{hackrf\_transfer}, GNU Radio File Sink with metadata logging.
  \end{itemize}

  \item \textbf{Latency Benchmark}
  \begin{itemize}
    \item Measure RF-to-processing pipeline delay under typical load.
    \item Use loopback test (TX $\rightarrow$ RX path with known signal).
    \item \textit{Pass criterion:} Latency meets real-time requirements (e.g., $<10$ ms for interactive applications).
    \item \textit{Tools:} GNU Radio \texttt{message\_strobe} with timestamped logging.
  \end{itemize}

  \item \textbf{Thermal Stress Test}
  \begin{itemize}
    \item Run at 100\% CPU + SDR acquisition for 4 hours minimum.
    \item Monitor core temperatures, clock speeds (throttling detection).
    \item \textit{Pass criterion:} No thermal shutdown, $<10\%$ performance degradation.
    \item \textit{Tools:} \texttt{stress-ng}, \texttt{sensors}, \texttt{turbostat}.
  \end{itemize}

  \item \textbf{Multi-Channel Verification}
  \begin{itemize}
    \item If applicable, run target number of simultaneous channels.
    \item Verify no cross-talk, interference, or resource contention.
    \item \textit{Pass criterion:} All channels maintain specifications simultaneously.
    \item \textit{Tools:} GNU Radio multi-USRP source, parallel \texttt{rtl\_sdr} instances.
  \end{itemize}

  \item \textbf{Synchronization Accuracy} (if multi-node or coherent operation)
  \begin{itemize}
    \item Verify timing alignment between nodes or channels.
    \item Measure phase drift over time with external reference.
    \item \textit{Pass criterion:} Timing jitter within application tolerance (e.g., $<100$ ns for beamforming).
    \item \textit{Tools:} PPS measurement, GNU Radio correlate blocks, external rubidium reference.
  \end{itemize}

  \item \textbf{Software Stack Integration}
  \begin{itemize}
    \item Install complete SDR toolchain (GNU Radio, drivers, dependencies).
    \item Verify version compatibility matrix (e.g., GNU Radio 3.10 + UHD 4.6).
    \item Run reference flowgraphs (FM receiver, waterfall, etc.).
    \item \textit{Pass criterion:} Clean installation, all blocks functional, no library conflicts.
    \item \textit{Tools:} Package manager logs, GNU Radio Companion examples.
  \end{itemize}

  \item \textbf{Container Deployment} (if applicable)
  \begin{itemize}
    \item Deploy SDR stack in Docker/Podman container.
    \item Verify USB/PCIe passthrough, GPU access (for ML).
    \item \textit{Pass criterion:} Container performance matches bare-metal within 5\%.
    \item \textit{Tools:} Docker Compose, device flags, GPU runtime.
  \end{itemize}

  \item \textbf{RF Performance Validation}
  \begin{itemize}
    \item Measure actual noise figure and sensitivity with known test equipment.
    \item Verify frequency accuracy and stability over temperature.
    \item Check for spurious emissions, images, harmonics.
    \item \textit{Pass criterion:} Meets datasheet specifications within tolerances.
    \item \textit{Tools:} Signal generator, spectrum analyzer, shielded enclosure.
  \end{itemize}
\end{enumerate}

\subsection*{Scored Evaluation Criteria}

\begingroup
\small

\criteriaheading{Compute Capability}
\begin{itemize}
  \scoreitem{5}{Exceptional}{High core count or equivalent parallel throughput, dedicated GPU/AI accelerator, enterprise-grade RAM; sustains 4+ simultaneous wideband channels with ML inference in real time.}
  \scoreitem{4}{Excellent}{Mid-range multicore or equivalent, integrated AI accelerator or GPU, substantial RAM; processes 2--3 channels with moderate ML workloads.}
  \scoreitem{3}{Adequate}{Standard multicore or equivalent, adequate RAM; handles single wideband channel or multiple narrowband channels.}
  \scoreitem{2}{Limited}{Entry-level core count, constrained memory; struggles with real-time processing, requires downsampling or offline analysis.}
  \scoreitem{1}{Insufficient}{Minimal processing capability, severely constrained memory; cannot sustain real-time baseband processing.}
\end{itemize}

\criteriaheading{SDR Software Compatibility}
\begin{itemize}
  \scoreitem{5}{Comprehensive}{Full SoapySDR module support, mature Python/C++ APIs with stable bindings, pre-built binaries for major OS/distros (Linux, Windows, macOS, ARM); works out of the box with GNU Radio/UHD/gr-osmosdr; active driver maintenance.}
  \scoreitem{4}{Strong}{SoapySDR module available (may require minor build), well-documented Python/C++ APIs, binaries available for Linux/ARM; community-maintained packages, comprehensive documentation.}
  \scoreitem{3}{Functional}{Basic SoapySDR support or limited API maturity, requires manual compilation for some OS/architectures, partial cross-platform binary availability; workarounds documented.}
  \scoreitem{2}{Partial}{No official SoapySDR module, immature or proprietary APIs, single-platform binaries only; significant porting effort, missing drivers for some SDR models.}
  \scoreitem{1}{Incompatible}{No SDR abstraction layer support, no usable APIs, architecture/OS locked, or proprietary stack only.}
\end{itemize}

\criteriaheading{I/O and RF Interface Support}
\begin{itemize}
  \scoreitem{5}{Comprehensive}{Multiple PCIe Gen3+ slots, several USB 3.x/4 ports (dedicated controllers), 10GbE+, Thunderbolt; sufficient bandwidth for high-throughput SDR arrays.}
  \scoreitem{4}{Strong}{PCIe Gen2+, multiple USB 3.x ports, GbE/2.5GbE; adequate bandwidth for 2+ SDRs without contention.}
  \scoreitem{3}{Adequate}{Dual USB 3.x (may share controller), GbE; handles single wideband SDR or light multi-channel setups.}
  \scoreitem{2}{Limited}{Shared USB bandwidth, may experience overflows at high sample rates without careful tuning.}
  \scoreitem{1}{Insufficient}{Legacy USB 2.0 or severely bandwidth-constrained interfaces; unsuitable for modern wideband SDR.}
\end{itemize}

\criteriaheading{Power and Thermal Characteristics}
\begin{itemize}
  \scoreitem{5}{Excellent}{Ultra-low power envelope, passive cooling capable, wide operating temperature range, zero throttling under sustained load.}
  \scoreitem{4}{Good}{Low-to-moderate power draw, efficient active cooling, minimal thermal throttling under sustained load.}
  \scoreitem{3}{Acceptable}{Moderate power envelope, requires active cooling, occasional thermal throttling under peak workloads.}
  \scoreitem{2}{Challenging}{High power draw, aggressive throttling without robust cooling infrastructure.}
  \scoreitem{1}{Problematic}{Very high power/thermal output or severe management issues; unsuitable for compact/field deployment.}
\end{itemize}

\criteriaheading{Form Factor and Ruggedness}
\begin{itemize}
  \scoreitem{5}{Field-Ready}{Fanless, high ingress protection rating, shock/vibration resistant, lightweight, compact footprint.}
  \scoreitem{4}{Portable}{Small form factor, ruggedized enclosure options, moderate weight, efficient fan-cooled design.}
  \scoreitem{3}{Desktop}{Mini-ITX or NUC-sized, standard PC enclosure, moderate weight, requires controlled environment.}
  \scoreitem{2}{Large}{Standard desktop/tower footprint, heavier, rack-mount or floor-standing, requires infrastructure.}
  \scoreitem{1}{Stationary}{Server-class chassis, heavy, requires dedicated data center or lab environment.}
\end{itemize}

\criteriaheading{Software Ecosystem and Updates}
\begin{itemize}
  \scoreitem{5}{Mature}{LTS Linux support, official container images, long-term update commitment, documented CI/CD pipelines.}
  \scoreitem{4}{Strong}{Mainline kernel support, community containers, multi-year updates, comprehensive documentation.}
  \scoreitem{3}{Functional}{Requires specific distro/kernel, community-maintained containers, periodic updates, adequate docs.}
  \scoreitem{2}{Fragile}{Requires out-of-tree patches, limited containerization, uncertain update cadence.}
  \scoreitem{1}{Problematic}{Abandonware risk, no container support, requires legacy software stacks.}
\end{itemize}

\criteriaheading{Expandability and Modularity}
\begin{itemize}
  \scoreitem{5}{Highly Modular}{Multiple expansion slots, RAM upgradeable to high capacity, NVMe/PCIe expansion, accelerator-ready interfaces.}
  \scoreitem{4}{Expandable}{RAM upgradeable to moderate capacity, 1--2 expansion slots, storage upgrades supported.}
  \scoreitem{3}{Limited Expansion}{Fixed base RAM, storage upgradeable, USB-centric expansion.}
  \scoreitem{2}{Minimal}{Soldered/limited RAM, eMMC/basic storage, dongle-only expansion.}
  \scoreitem{1}{Fixed}{No upgrade path, proprietary form factor or sealed architecture.}
\end{itemize}

\criteriaheading{Interoperability and Openness}
\begin{itemize}
  \scoreitem{5}{Fully Open}{Open-source drivers in mainline, fully documented APIs, public reference designs, zero vendor lock-in.}
  \scoreitem{4}{Mostly Open}{Open drivers (may require build), well-documented APIs, community reverse-engineering support if needed.}
  \scoreitem{3}{Mixed}{Binary blobs required, partial documentation, some vendor tools necessary for full functionality.}
  \scoreitem{2}{Proprietary}{Closed drivers, restricted documentation, vendor SDK mandatory for development.}
  \scoreitem{1}{Locked}{Encrypted/undocumented firmware, no alternative stacks possible, hard vendor lock-in.}
\end{itemize}

\criteriaheading{Reliability and Lifecycle}
\begin{itemize}
  \scoreitem{5}{Industrial}{High MTBF rating, extended warranty, decade-plus lifecycle commitment, guaranteed long-term part availability.}
  \scoreitem{4}{Commercial}{Solid MTBF, multi-year warranty, long lifecycle, established supply chain.}
  \scoreitem{3}{Standard}{Typical MTBF, standard warranty, multi-year lifecycle, standard commercial support.}
  \scoreitem{2}{Consumer}{Lower MTBF, limited warranty, short lifecycle, typical consumer-grade support.}
  \scoreitem{1}{Short-Lived}{No formal MTBF, minimal warranty, high EOL risk, unsuitable for sustained operations.}
\end{itemize}

\criteriaheading{Total Cost of Ownership}
\begin{itemize}
  \scoreitem{5}{Very Low}{Minimal upfront cost, ultra-low power draw, zero licensing, predictable long-term operating expenses.}
  \scoreitem{4}{Low}{Moderate upfront cost, low power draw, minimal licensing, manageable long-term expenses.}
  \scoreitem{3}{Moderate}{Mid-range upfront cost, moderate power/licensing, standard operational overhead.}
  \scoreitem{2}{High}{Significant upfront investment, high power/licensing costs, substantial maintenance overhead.}
  \scoreitem{1}{Very High}{Premium upfront cost, excessive power/licensing, expensive support/upgrade contracts.}
\end{itemize}

\criteriaheading{Multi-Node Readiness}
\begin{itemize}
  \scoreitem{5}{Orchestration-Ready}{Native cluster/orchestration support, centralized management APIs, distributed precision timing, validated at scale.}
  \scoreitem{4}{Cluster-Capable}{Supports orchestration with configuration, centralized monitoring, external timing reference, validated for mid-size deployments.}
  \scoreitem{3}{Multi-Node}{Manual cluster setup, SSH-based management, basic GPS sync capable, suitable for small deployments.}
  \scoreitem{2}{Networked}{Basic network connectivity, independent node management, no hardware timing sync, limited scalability.}
  \scoreitem{1}{Standalone}{Single-node architecture, no clustering or distributed management capability.}
\end{itemize}

\criteriaheading{Security and Manageability}
\begin{itemize}
  \scoreitem{5}{Hardened}{Hardware-backed secure boot, TPM/HSM support, hardware-encrypted storage, enterprise remote management, FIPS-compliant options.}
  \scoreitem{4}{Secure}{Secure boot, software encryption, remote management via standard protocols, RBAC, centralized logging.}
  \scoreitem{3}{Managed}{Basic secure boot, encrypted partitions, standard remote access, local logging.}
  \scoreitem{2}{Basic}{Manual update process, password-only auth, no default encryption.}
  \scoreitem{1}{Insecure}{No secure boot, no remote management, no encryption options, exposed attack surface.}
\end{itemize}

\criteriaheading{RF Fidelity and Sensitivity}
\begin{itemize}
  \scoreitem{5}{Laboratory-Grade}{Ultra-low noise figure, exceptional dynamic range, minimal phase noise, factory-calibrated reference oscillator.}
  \scoreitem{4}{Professional}{Low noise figure, high dynamic range, stable phase noise profile, precision reference (TCXO/GPSDO).}
  \scoreitem{3}{Good}{Moderate noise figure, standard dynamic range, acceptable phase noise, external reference capable.}
  \scoreitem{2}{Adequate}{Higher noise floor, limited dynamic range, noticeable phase noise, basic timing reference.}
  \scoreitem{1}{Poor}{High noise figure, severely limited dynamic range, unstable oscillator, unsuitable for precision/weak-signal work.}
\end{itemize}

\criteriaheading{Timing and Synchronization}
\begin{itemize}
  \scoreitem{5}{Coherent Array}{Sub-nanosecond sync capability, integrated precision timing, dedicated 10 MHz/1PPS I/O, multi-channel phase-locked.}
  \scoreitem{4}{Precision Sync}{High-precision network sync, external precision timing support, 10 MHz/1PPS input, dual-channel coherence.}
  \scoreitem{3}{Synchronized}{External reference input, software-level network sync, basic GPS pulse support.}
  \scoreitem{2}{Basic Timing}{Standard network time only, no external reference, independent free-running clocks.}
  \scoreitem{1}{Unsynchronized}{No timing infrastructure, oscillator drift unmanaged, unsuitable for multi-node/coherent work.}
\end{itemize}

\criteriaheading{FPGA/DSP Offloading}
\begin{itemize}
  \scoreitem{5}{Open HDL}{User-accessible FPGA with open/free toolchains, documented host interfaces, comprehensive example designs.}
  \scoreitem{4}{Programmable}{Accessible FPGA with vendor tools (free tier sufficient), documented API for host interaction, community bitstreams.}
  \scoreitem{3}{Semi-Custom}{Fixed base bitstream with configurable blocks, vendor SDK required for modifications.}
  \scoreitem{2}{Fixed Function}{Closed FPGA architecture, firmware-only updates, minimal customization options.}
  \scoreitem{1}{No Access}{Black-box processing or no FPGA/DSP offload available; fully host-dependent.}
\end{itemize}

\criteriaheading{Regulatory and Compliance}
\begin{itemize}
  \scoreitem{5}{Fully Certified}{Complete regulatory certification for enclosure, emissions compliant, integrated shielding/filtering, documented compliance reports.}
  \scoreitem{4}{Compliant}{Enclosure meets standard EMI/EMC guidelines, external filtering available, vendor compliance documentation.}
  \scoreitem{3}{Adequate}{Basic internal shielding, user-responsible for final compliance, commercial-grade enclosure.}
  \scoreitem{2}{Limited}{Minimal shielding, may require custom enclosure for certification, suitable for lab/hobby use only.}
  \scoreitem{1}{Non-Compliant}{No shielding, known emission issues, unsuitable for regulated or production environments.}
\end{itemize}

\endgroup

\subsection*{3.1 Knock-Out Thresholds}

\subsubsection*{Per-Template Minimum Viable Scores}

\begin{table}[H]
  \centering
  \small
  \begin{tabular}{@{}p{0.29\textwidth}p{0.63\textwidth}@{}}
    \toprule
    Template & Knock-Out Criteria (Auto-Reject if any score is below threshold) \\
    \midrule
    1. Spectrum Monitoring \& SIGINT & Compute capability $\geq 3$; Security / manageability $\geq 3$; RF fidelity $\geq 3$; Timing \& synchronization $\geq 3$; Multi-node readiness $\geq 3$ \\
    2. Field Research \& Portable Ops & Compute capability $\geq 3$; Power / thermal $\geq 3$; Form factor $\geq 3$; Reliability $\geq 3$ \\
    3. Laboratory Research \& Algorithm Dev & Compute capability $\geq 3$; SDR software compatibility $\geq 3$; I/O interfaces $\geq 3$; Software ecosystem $\geq 3$; Openness $\geq 3$ \\
    4. Production IoT / Wireless Monitoring & Reliability $\geq 3$; Total cost of ownership $\geq 3$; Multi-node readiness $\geq 3$; Security / manageability $\geq 3$; Regulatory compliance $\geq 3$ \\
    \bottomrule
  \end{tabular}
  \caption{Minimum viable scores by deployment template. A platform failing any listed criterion is non-viable for that role.}
  \label{tab:knockout-thresholds}
\end{table}

\subsubsection*{Rationale for Minimums}

\begin{itemize}
  \item \textbf{Template 1 (SIGINT):} Continuous, multi-site coherent sensing cannot tolerate insufficient compute (dropped samples), weak security (exposure of classified/controlled data), poor RF fidelity (missed weak signals), or broken timing (loss of coherence).
  \item \textbf{Template 2 (Field):} Battery-powered, single-operator deployments fail outright if thermal design forces shutdown, the chassis cannot survive transport, or compute cannot sustain real-time acquisition.
  \item \textbf{Template 3 (Lab):} Experimentation requires open toolchains, stable APIs, and sufficient I/O to iterate. A closed or fragile ecosystem blocks research progress regardless of raw performance.
  \item \textbf{Template 4 (Production IoT):} Fleet deployments prioritize unit economics, compliance, and remote manageability. High TCO or missing regulatory certification makes scale impossible.
\end{itemize}

\clearpage

\subsubsection*{Decision Matrix Output Template}

Use this 1-page summary for final platform selection:

\decisionsummarybox

\subsubsection*{Scoring Table Integration}

Replace the final rows of your platform-specific scoring table with:

{\small
\begin{verbatim}
\midrule
\multicolumn{2}{l}{\textbf{Raw Weighted Total:}}
& \underline{\hspace{2cm}} / \underline{\hspace{2cm}} \\
\multicolumn{2}{l}{\textbf{Applicable Criteria (N):}}
& \underline{\hspace{2cm}} of 16 \\
\multicolumn{2}{l}{\textbf{Normalized Score (\%):}}
& \underline{\hspace{2cm}} \% \\
\multicolumn{2}{l}{\textbf{Decision (Tab.~\ref{tab:decision-thresholds}):}}
& \underline{\hspace{4cm}} \\
\multicolumn{2}{l}{\textbf{Knock-Out Violations (Tab.~\ref{tab:knockout-thresholds}):}}
& \underline{\hspace{4cm}} \\
\midrule
\textit{Pros:} & \multicolumn{2}{l}{} \\
\textit{Cons/Caveats (incl. N/A rationale):} & \multicolumn{2}{l}{} \\
\textit{Reviewer / Date:} & \multicolumn{2}{l}{\today} \\
\bottomrule
\end{verbatim}
}

\end{document}
